\subsection{All finite changes of a nonzero integer at the finite level-three quotient}

\begin{statement}{A.68}{Exact finite-alteration exponent sets at the actual level-three quotient}{Proved}
On $E:y^2+y=x^3-1$, $P=(1,0)$, with the
punctures $E[3]$ and the actual joint quotient
$\Gamma_3$ of A.67, the following exact
statements hold.
\begin{enumerate}
\item If $a\in\widehat\Z$ agrees with a
      nonzero integer $n$ outside a finite
      set of coefficient primes, then
      $a\kappa(P)$ has a $\Gamma_3$ Selmer
      lift if and only if $a=n$.
\item If $a$ agrees with any integer $n$,
      including zero, outside coefficient
      prime three, then it has such a
      lift if and only if $a=n$.
\end{enumerate}
For every elliptic cubic twist with the
rational infinity flex and notation of A.67,
the first exclusion holds whenever at least
one changed coefficient prime is different
from three. Thus the assertion includes
the coefficient prime two, and simultaneous
changes at finitely many primes.

These are statements about the finite
joint congruence quotient, stronger on
these exponent families than merely
excluding a full abelian Selmer lift.
They do not determine the finite quotient's
fibres or arbitrary infinite alterations.
A.69 completes the finite-support exponents
around zero with other coefficient primes.
\end{statement}
\begin{proof}
\par\smallskip\noindent\textbf{Prescribing a primary component.}\enspace 
We first make explicit the strengthening
of A.61 used here. Let $F$ contain the
CM endomorphisms, let $\mathcal P$ be a
finite coefficient-prime set, and let
$R\in E(F)$ be nontorsion. After a finite
extension, prescribe any rational
$p$-primary torsion target $T_p$ for
each $p\in\mathcal P$.
The joint Kummer image is open in the
product of the groups
$T_pE\rtimes G_p$, by A.61. It therefore
contains a product neighbourhood
\[
 \prod_{p\in\mathcal P}
        (p^{b_p}T_pE\times G_{p,1})
 \tag{A68.1}
\]
of the identity, with $G_{p,1}$ open.
Lift $T_p$ to $\tau_p\in V_pE$.
Choose $A_p\in G_{p,1}$ sufficiently
close to the identity, with $A_p-1$
invertible, so that
$c_p=(A_p-1)\tau_p\in p^{b_p}T_pE$.
Such matrices occur arbitrarily close
to the identity: sufficiently high
$p$-powers of a good Frobenius, after
killing its finite prime-to-$p$ part,
converge to the identity, and no
positive power has eigenvalue one
because its eigenvalues have nonzero
Weil weight. For the zero target
take $\tau_p=0$.

The tuple $(A_p,c_p)_p$ lies in
the actual joint image. Near it,
the conditions
\[
 c_p-(A_p-1)\tau_p\in(A_p-1)T_pE
                   \qquad(p\in\mathcal P)
 \tag{A68.2}
\]
are open. Indeed the denominator
valuation of $(A_p-1)^{-1}$ is
locally bounded and the image
lattice is constant on a sufficiently
small congruence neighbourhood.
At a good place whose residue
characteristic is outside
$\mathcal P$, (A68.2) is equality
of the unramified local Kummer
classes of $R$ and $T_p$.
Full local Kummer injectivity
therefore identifies the actual
$p$-primary component of $R$
with $T_p$. The condition is
conjugacy invariant when the
targets are rational: translation
conjugacy changes the cocycle by
$(A_p-1)T_pE$. We have obtained
a nonempty open Frobenius condition
prescribing all the primary targets
simultaneously. This argument is
reapplied after every finite base
extension below.

\par\smallskip\noindent\textbf{A changed prime different from three.}\enspace 
Let $a_q=n$ outside a finite set
$\mathcal P$, where $n\ne0$, and
put $\delta_p=a_p-n$. Suppose
$p\in\mathcal P$, $p\ne3$, and
$\delta_p\ne0$. Choose one flex-ratio
unit $g$ of A.67 over
$L=\Q(E[3])$, with
\[
 \operatorname{div}(g)=3(d_1)-3(d_0),
              \qquad d_0,d_1\in E[3](L),\ d_0\ne d_1.
 \tag{A68.3}
\]
Its quadratic Kummer character
factors through $\Gamma_{3,L}$.
A hypothesized Selmer lift supplies
one global class $\beta\in
H^1(L,\mu_2)$ with the corresponding
local point evaluations.

Fix once a finite extension $F/L$
containing the CM endomorphisms,
every $E[\ell]$ for $\ell\in\mathcal P$,
and a square root of a representative
of $\beta$. The class $\beta$
is zero over $F$. Let $K_\infty/F$
be the full joint torsion and
$R$-division tower for the primes
in $\mathcal P$.
There is a single finite set $S$
outside which this entire tower
is unramified. Include $2$,
$\mathcal P$, bad reduction and
the fixed field's ramification
in $S$. At every other place,
$E$ is an abelian scheme and
multiplication by every
$\ell^r$ is finite \'{e}tale.
The fibre above $R$ is finite
\'{e}tale as well: properness
makes $R$ an integral section.
This proves uniform unramifiedness
of all the torsion and division
extensions, not just a bounded
finite layer.

For an arbitrary $p$-primary
torsion point $T$, choose a
$p$-primary $T'$ with
$\delta_pT'=T$. Such a point
exists by dividing by the finite
$p$-adic valuation of $\delta_p$
and inverting its unit part.
Take a finite torsion field
$F_r\subset K_\infty$ defining
both points. The value
$h_T=g(nR+T)$ is nonzero and
finite, since $nR+T$ is
nontorsion and (A68.3) is
supported on torsion.

We claim that $\sqrt{h_T}\in
K_\infty$. Otherwise the
quadratic extension
$F_r(\sqrt{h_T})$ has trivial
intersection with the Galois
tower $K_\infty/F_r$.
Chebotarev can then impose a
nontrivial quadratic Frobenius
and the open condition above
with $R_p=T'$ and every other
altered primary component zero.
At a resulting good place the
actual compact local centre is
\[
 aR=nR+\sum_{\ell\in\mathcal P}
                    \delta_\ell R_\ell
                  =nR+T.
 \tag{A68.4}
\]
It is interior. The zero global
class $\beta$ forces $g(nR+T)$
to be a local square, contrary
to the chosen Frobenius.
This proves the claim, including
the case where the radical was
already in $F_r$.

Consequently, for every $T$,
the quadratic extension defined
by $g(nR+T)$ is unramified
outside the same fixed $S$.
Its valuations there are even,
since the residue characteristic
is not two. The evenness holds
in any finite torsion field
defining $T$: the larger fields
used to define $T'$ are
unramified outside $S$, so their
ramification indices there are
one. This is valid also when
$p=2$; residue degrees play
no role in this valuation
descent.

Set $m=p^2$. Then $m\equiv1
\pmod3$, so $[m]$ fixes $d_0,d_1$.
Pushing the divisor (A68.3)
forward along $[m]$ gives
\[
 \operatorname{Norm}_{[m]}g=Cg,
                      \qquad C\in F^\times.
 \tag{A68.5}
\]
Iterating the norm, whose degree
is $m^2$, gives
$C_j=C^{\,1+m^2+\cdots+m^{2j-2}}$
and the exact fibre identity
\[
 \prod_{T\in E[m^j]}g(nR+T)
                  =C_j\,g([m^j]nR).
 \tag{A68.6}
\]
Enlarge $S$ by the support of
$C$. In a finite torsion field
containing $E[m^j]$, every
factor on the left has even
valuation outside $S$.
So does the right-hand value
$g([m^j]nR)$. This value lies
in $F$, and the extension is
unramified there, so the same
evenness holds in $F$ itself.

The infinitely many distinct
points $[m^j]nR$ therefore lift
to the finitely many twists
$z^2=c^{-1}g$, with
$c\in F(S,2)$. Finiteness of
$F(S,2)$ follows from the
$S$-unit theorem and the
finite ideal-class group.
Their compactifications have
genus two by (A68.3).
Faltings's theorem contradicts
infinitely many such points
\citep[Satz~7 and the correction]{AMFaltings1983,AMFaltings1984Erratum}.
This proves the uniform
exclusion whenever a changed
prime is different from three.
In particular it proves it
for two and for simultaneous
finite changes. At no point
is a quadratic extension
assumed disjoint from a
pro-two tower when it is
contained in that tower;
containment is precisely
what supplies the uniform
ramification bound.

\par\smallskip\noindent\textbf{Arbitrary changes at coefficient three.}\enspace 
Now specialize to the least
curve and write $a_3=\alpha$,
$a_q=n$ for $q\ne3$.
Use $L$, $d_1$, and $g$ of
(A67.12)--(A67.14).
Over $L$ the proper
$3$-torsion and $P$-Kummer
tower is pro-$3$.
The zero-primary Frobenius
condition and quadratic
disjointness, exactly as
in A.67, force the global
$g$-class to be
$\kappa_2(g(nP))$ when
$n\ne0$. When $n=0$, they
force the zero class:
the local centre at those
places is the cusp $O$,
but $g$ extends as a regular
unit at $O$ with $g(O)=1$.
Its value is independent
of every tangent parameter,
so cusp inertia contributes
nothing.

Suppose first that $n\ne0$
and $\alpha\ne n$. Put
$k=v_3(\alpha-n)$ and
choose $T\in E[3^{k+1}]$
with $(\alpha-n)T=d_1$.
Over $L'=L(E[3^{k+1}])$
this point is rational.
The extension $L'/L$
has odd $3$-power degree.
By (A67.14),
\[
 h=\frac{g(nP+d_1)}{g(nP)}
 \tag{A68.7}
\]
is nonsquare in $L$, and
it stays nonsquare over
$L'$. Its quadratic extension
is disjoint from the pro-$3$
tower over $L'$.
The prescribed-primary lemma
now permits both a nonsquare
quadratic Frobenius and
$P_3=T$. At this place the
forced compact centre is
$nP+(\alpha-n)T=nP+d_1$,
contradicting the forced
global class $g(nP)$.
Therefore $\alpha=n$.

Finally let $n=0$ and
$\alpha\ne0$. We exhibit
a fixed order-nine point
$S$ for which $g(S)$ is
nonsquare in $L_9=L(E[9])$.
At the good split prime
$73$, take $b\equiv25$,
$\zeta\equiv8$. Exact
addition gives
\[
 \operatorname{ord}_{E(\mathbf F_{73})}
           (1,0)=9,\qquad
 g(1,0)=56,\qquad 56^{36}=-1.
 \tag{A68.8}
\]
Lift this prime-to-$73$
torsion point uniquely to
$S\in E[9](\Q_{73})$.
It is an algebraic point
defined over $L_9$.
Every residue degree in
the pro-$3$ extension
$L_9/L$ is odd, so the
nonsquare residue in
(A68.8) persists at the
chosen place of $L_9$.
Thus $g(S)$ is globally
nonsquare. Its value is an
ordinary interior evaluation:
$S$ has exact order nine,
so $S\notin E[3]$.

Choose $T$ with $\alpha T=S$
over $L_r=L(E[3^{v_3(\alpha)+2}])$.
The odd extension preserves
the nonsquare class of
$g(S)$. Its quadratic
extension is disjoint from
the pro-$3$ tower over
$L_r$. Prescribed-primary
Chebotarev with $P_3=T$
therefore gives the local
centre $aP=\alpha T=S$
and a nonsquare value of
$g(S)$. This contradicts
the zero global $g$-class.
Hence $\alpha=0$.

Every integer exponent is
realized by its point
section when nonzero and
by a cusp section at $O$
when zero. Combining the
prime-three argument with
the norm argument proves
both exact exponent-set
claims.
\end{proof}

The data (A68.8) are already
present in the executed
\texttt{a\_main\_jline\_at3.out}
certificate. The torsion-point
and norm arguments above are
proved uniformly and do not
depend on extrapolating a
finite calculation. The
nonzero-baseline hypothesis
in the norm part is essential
to its stated proof: when
$n=0$, the points $[m^j]nR$
are constant, so its Faltings
contradiction disappears.
Only the separately proved
pure coefficient-three
case includes that baseline.
